\chapter*{Introduction générale}
\addcontentsline{toc}{chapter}{Introduction Générale}


La transformation numérique est une opportunité stratégique dans laquelle le secteur des services publics a la possibilité de se développer afin d'optimiser sa performance et de renforcer les relations administratifs citoyens. En effet en raison des évolutions rapides des technologies de l'information et de la communication les services publics doivent s adapter pour être plus efficaces et proches de leurs usagers.

En République Centrafricaine RCA la gestion des services administratifs est actuellement assurée de manière décentralisée par plusieurs administrations chacune disposant de ses propres procédures et outils. Cette organisation bien que fonctionnelle engendre certaines limites telles qu un manque de visibilité globale sur les services offerts des difficultés dans le suivi des demandes citoyennes des délais de traitement importants ainsi qu une communication parfois insuffisante entre les institutions publiques et les citoyens.

Face à ces constats le ministère des finances et du budget de la République Centrafricaine a exprimé le besoin de mettre en place une plateforme numérique centralisée permettant de fédérer structurer et piloter l'ensemble des services administratifs du pays. Cette plateforme vise à offrir un cadre unifié sécurisé et évolutif pour la gestion des administrations des demandes des réclamations et des échanges entre les différents intervenants du système.

Dans ce contexte ce projet de fin d études vise à concevoir et développer une application web intégrée destinée à accompagner la transformation numérique des services publics en assurant leur digitalisation. La solution proposée permet aux ministères de gérer leurs services de manière dynamique aux responsables administratifs d'assurer le suivi et le traitement des demandes et réclamations et aux citoyens de déposer suivre et évaluer leurs démarches administratives en ligne en toute transparence. Elle intègre également un assistant conversationnel intelligent qui guide les citoyens et répond à leurs questions en langage naturel.

Ce rapport est organisé en quatre chapitres structurés comme suit.

\begin{itemize}

    \item \textbf{Chapitre 1 — Etude préalable :} comprend la description de l entité accueillante et du contexte dans lequel est réalisé le projet On y explique la problématique fait une analyse des solutions actuelles par benchmarking de différentes plateformes et on expose la solution retenue et les les méthodes de gestion de projet mises en œuvre.

    \item \textbf{Chapitre 2 — Spécification des besoins :} concerne l identification des différents acteurs du système la détermination des exigences fonctionnelles et non fonctionnelles la création du backlog du produit l organisation des sprints distribués sur deux releases et l exposition de l architecture choisie et des technologies utilisées.

    \item \textbf{Chapitre 3 — Release 1 :} couvre les \textit{Sprint 1} \textit{Sprint 2} et \textit{Sprint 3} Il détaille la mise en place de la gouvernance des accès et des entités organisationnelles la configuration du catalogue de services administratifs avec le moteur de formulaires dynamiques et la dématérialisation des demandes des citoyens.

    \item \textbf{Chapitre 4 —Release 2 :} couvre les \textit{Sprint 4} et \textit{Sprint 5} Il présente le système de gestion des réclamations le moteur de notifications et les outils de supervision de la plateforme puis le portail public avec la vitrine institutionnelle et un assistant conversationnel intelligent fondé sur l'intelligence artificielle approche RAG et son canal d'assistance citoyen.

\end{itemize}

Le rapport se termine par une conclusion générale qui synthétise les principaux résultats obtenus lors du projet ainsi que les améliorations possibles et futures.